\session{2}{10/10}{1限}{1}{情報と意思決定}

\goals{
  \item データ・情報・知識・知恵の階層（DIKW）を例を挙げて説明できる。
  \item サイモンの意思決定プロセス（情報活動・設計活動・選択活動）と限定合理性を理解する。
  \item 組織の階層によって必要な情報の性質が異なることを理解し、AIが各段階にどう関わるかを考える。
}

\guidesec{解説}

\topic{データ・情報・知識の階層}
「情報」という言葉は日常では広く使われますが、経営情報では次のように段階を分けて考えます（DIKW 階層）。
\par\noindent\begin{gtabh}{colspec={Q[l,wd=18mm] X[l] X[l]}}
  段階 & 定義 & 例（小売店） \\
  データ & 事実や出来事の記録。それだけでは意味を持たない。 & 「10月9日、商品Aが12個売れた」 \\
  情報 & データに文脈や意味を与え、受け手にとって意味を持つようにしたもの。 & 「商品Aの売上は前週比で30\%増えた」 \\
  知識 & 情報を体系化し、判断や行動に使えるようにしたもの。 & 「気温が上がると商品Aが売れるので、暑い日は仕入れを増やす」 \\
  知恵 & 知識を状況に応じて使い分け、新しい状況にも対応する力。 & 「今年は猛暑でも競合が値下げしたので、仕入れ増は見送る」 \\
\end{gtabh}
生成AIは、データから情報への変換（集計・要約）と、情報から知識への変換（パターンの発見・一般化）を速く行えます。
一方で、知恵の段階、つまり自社の状況や価値観に照らした使い分けは、人が担う部分が大きいと言えます。

\topic{サイモンの意思決定プロセス}
ハーバート・サイモン（1978年ノーベル経済学賞）は、意思決定を次の3段階に分けました。
\begin{itemize}
  \item \textbf{情報活動（Intelligence）}　問題や機会を発見し、意思決定が必要な状況を認識する。環境を探索し、データを集める。
  \item \textbf{設計活動（Design）}　問題に対する解決策の代替案を考え出し、それぞれの結果を予測する。
  \item \textbf{選択活動（Choice）}　代替案の中から1つを選ぶ。
\end{itemize}
後にこれに \textbf{再検討（Review）}、つまり実行した結果を評価して次の意思決定に活かす段階が加えられました。
実際の意思決定はこの順に一直線に進むのではなく、設計活動の途中で情報活動に戻るなど、行き来しながら進みます。

生成AIはどの段階にも関わりますが、特に相性が良いのは設計活動です。人が思いつく代替案は数個ですが、
AIは数十案を短時間で出し、それぞれの長所・短所を整理できます。情報活動では情報収集と要約、選択活動では
評価基準の整理と比較表の作成に使えます。ただし、最終的に選ぶのは人です。

\topic{限定合理性と満足化}
古典的な経済学は、意思決定者がすべての代替案とその結果を知っていて、最適なものを選ぶと仮定しました（完全合理性）。
サイモンはこれを批判し、人の情報収集・処理能力には限界があるため、すべてを検討することはできず、
「これで十分」と思える水準（満足水準）を満たす案が見つかった時点で選択を終える、と考えました。
これが \textbf{限定合理性（bounded rationality）} と \textbf{満足化（satisficing）} です。

AIは人より多くの代替案を検討でき、限定合理性の限界を広げます。しかし、AIにも学習データや設計に由来する限界があり、
「AIが出した案だから最適だ」とは言えません。むしろ、AIが示す候補の中から「何を満足水準とするか」を決めるのが人の役割になります。

\topic{組織の階層と情報の流れ}
アンソニーは経営活動を3つの階層に分けました。階層によって、意思決定の性質と必要な情報が異なります。
\par\noindent\begin{gtabh}{colspec={Q[l,wd=28mm] X[l] X[l] X[l]}}
  階層 & 意思決定の例 & 意思決定の性質 & 必要な情報 \\
  戦略的計画（経営層） & 新規事業への参入、工場の立地 & 非定型、長期、不確実 & 要約された情報、外部情報、将来予測 \\
  マネジメントコントロール（管理層） & 部門の予算配分、人員計画 & 半定型、中期 & 部門別の実績と計画の比較 \\
  業務統制（現場） & 発注量の決定、シフト作成 & 定型、短期、反復 & 詳細で正確な内部情報、リアルタイム \\
\end{gtabh}
組織の中では、現場のデータが集約されて上位層に報告され（ボトムアップ）、上位層の方針が指示として現場に伝わります（トップダウン）。
情報システムは、この上下の流れと、部門間の横の流れを支えています。定型的な意思決定ほど自動化しやすく、
AIの導入もまず業務統制の階層から進んでいます。

\topic{キーワード}
\par\noindent\begin{keywords}{}
  DIKW階層 & データ・情報・知識・知恵の4段階。下から上へ、文脈と判断が加わる。 \\
  情報活動・設計活動・選択活動 & サイモンの意思決定の3段階。問題の発見、代替案の作成、選択。 \\
  限定合理性 & 人の情報処理能力には限界があり、最適解ではなく満足できる解を選ぶという考え方。 \\
  満足化 & 満足水準を満たす最初の案を選ぶ行動。最適化の対概念。 \\
  アンソニーの3階層 & 戦略的計画・マネジメントコントロール・業務統制。階層により情報の性質が異なる。 \\
  定型／非定型の意思決定 & 手順が決まっていて反復される意思決定と、その都度判断が必要な意思決定。 \\
\end{keywords}

\guidesec{演習課題}

\exercise{1}{AI演習：身近な意思決定をAIと対話しながら段階に分解する}
\instr{自分が最近行った意思決定（進路、アルバイト先の選択、大きな買い物など）を1つ選び、次のプロンプトでAIと対話しながら、サイモンの3段階に分解してください。分解したら、各段階で「実際に自分がしたこと」と「AIがいれば任せられたこと」を分けて書いてください。}
\begin{promptbox}[意思決定の分解]
私は大学生です。最近「（意思決定の内容）」を決めました。この意思決定を、サイモンの意思決定プロセス（情報活動・設計活動・選択活動）の3段階に分けて整理したいので、各段階について私に質問をしてください。一度に1つずつ質問し、私の回答を受けて次の質問をしてください。
\end{promptbox}

\exercise{2}{DIKW階層を自分の例で説明する}
\instr{大学の授業、スポーツ、SNS のいずれかを題材に、データ・情報・知識・知恵の4段階に対応する具体例を1つずつ書いてください。各段階で「前の段階に何が加わったか」も書きます。}

\exercise{3}{組織の階層と情報の性質}
\instr{コンビニエンスストアのチェーンを例に、「店舗の店員」「店舗を束ねるエリアマネジャー」「本社の経営層」が行う意思決定を1つずつ挙げ、それぞれに必要な情報の性質（定型か非定型か、詳細か要約か、内部か外部か、短期か長期か）を表にまとめてください。}

\guidesec{模範解答}

\begin{answer}{演習1}
例：アルバイト先の選択
\par\noindent\begin{gtabh}{colspec={Q[l,wd=20mm] X[l] X[l]}}
  段階 & 実際にしたこと & AIに任せられたこと \\
  情報活動 & 「来月から生活費が足りない」と気づいた。求人サイトで条件を見た。 & 求人情報の収集と、条件（時給・場所・時間帯）での絞り込み・一覧化。 \\
  設計活動 & 近所のカフェ、塾講師、家庭教師の3つを候補にした。 & 候補の洗い出し（自分では思いつかなかった大学内の仕事など）と、各候補の長所・短所の整理。 \\
  選択活動 & 授業との両立を重視して塾講師を選んだ。 & 評価基準ごとの比較表の作成。ただし「何を重視するか」は自分で決めた。 \\
  再検討 & 1か月後、授業の準備時間が減ったことに気づいた。 & 実績（シフト時間と学習時間）の記録と振り返りの質問。 \\
\end{gtabh}
AIと対話して気づく点：自分の意思決定は情報活動が不十分（候補を2〜3件しか見ていない）で、設計活動が省略されがち
（最初に見つけた案に決める＝満足化）だということ。AIを使うと代替案が増えるが、評価基準を自分で持たないと選べない。
\end{answer}

\begin{answer}{演習2}
例：SNS
\begin{itemize}
  \item データ：「10月8日の投稿に『いいね』が84件、閲覧が1,200件」。加わったもの：なし（事実の記録）。
  \item 情報：「写真付きの投稿は、文字だけの投稿より平均で閲覧が2倍多い」。加わったもの：比較という文脈。
  \item 知識：「新製品の紹介は写真付きで、夕方に投稿すると反応が良い」。加わったもの：複数の情報から一般化したパターン。
  \item 知恵：「今回は災害直後なので、宣伝の投稿は控える」。加わったもの：状況と価値観に照らした判断。
\end{itemize}
\end{answer}

\begin{answer}{演習3}
\par\noindent\begin{gtabh}{colspec={Q[l,wd=24mm] X[l] Q[l,wd=16mm] Q[l,wd=16mm] Q[l,wd=16mm] Q[l,wd=16mm]}}
  立場 & 意思決定の例 & 定型性 & 詳細さ & 情報源 & 時間軸 \\
  店員 & 廃棄になりそうな弁当の値引きをいつ行うか & 定型 & 詳細（商品ごと） & 内部（POS） & 短期（時間単位） \\
  エリアマネジャー & 担当店舗のうち、どの店の発注を見直すか & 半定型 & 店舗別の要約 & 内部中心、一部外部（周辺の競合） & 中期（週〜月） \\
  経営層 & 新しい業態（無人店舗）に投資するか & 非定型 & 高度に要約 & 外部（市場、技術動向）中心 & 長期（年単位） \\
\end{gtabh}
店員の意思決定は POS データに基づく定型的なもので、すでにAIによる値引き推奨が導入されている例があります。
経営層の意思決定は外部情報と将来予測に依存し、AIは情報収集と代替案の整理を助けますが、最終判断は人が行います。
\end{answer}

\guidesec{出席確認クイズの解説}
講義サイトの出席確認ページ（第2回）の5問です。
% 出席確認クイズ 第02回　情報と意思決定
%   attendance/attendance02.html から生成（解説は本ガイド用に加筆）。

\quizq{1}{データ・情報・知識の階層のうち、「データに文脈や意味を与えたもの」はどれでしょうか?}%
  {知恵}%
  {\correct{情報}}%
  {データ}%
  {雑音}%
  {正解は (b)。DIKW 階層では、データ（事実の記録）に文脈や意味を与えたものが情報、情報を体系化して活用できるようにしたものが知識、知識を状況に応じて使い分ける力が知恵です。たとえば「気温32度」はデータ、「例年より5度高い」は情報、「暑い日はアイスが売れるので仕入れを増やす」は知識にあたります。(d) の雑音はデータに混じる意味のない部分です。}

\quizq{2}{サイモンの意思決定プロセスの3つの段階に含まれ\underline{ない}ものはどれでしょうか?}%
  {選択活動(Choice)}%
  {\correct{宣伝活動}}%
  {設計活動(Design)}%
  {情報活動(Intelligence)}%
  {正解は (b)。サイモンは意思決定を情報活動（Intelligence）・設計活動（Design）・選択活動（Choice）の3段階で説明しました。後に再検討（Review）を加えることもあります。「宣伝活動」は含まれません。英語名と日本語名の対応を覚えておくと、文献を読むときに役立ちます。}

\quizq{3}{意思決定プロセスのうち、複数の代替案を作り出す段階はどれでしょうか?}%
  {選択活動(Choice)}%
  {\correct{設計活動(Design)}}%
  {情報活動(Intelligence)}%
  {評価活動}%
  {正解は (b)。設計活動は、問題に対する解決策の代替案を考え出す段階です。情報活動は問題や機会を発見する段階、選択活動は代替案の中から1つを選ぶ段階です。(d) の評価活動という名称は3段階にはありません。生成AIは代替案を数多く出せるため、設計活動との相性が特に良いと言えます。}

\quizq{4}{「限定合理性」の説明として適切なものはどれでしょうか?}%
  {\correct{人は情報処理能力の限界のもとで、満足できる水準の選択をする}}%
  {AIには合理性がまったくない}%
  {合理的な判断は経営には不要である}%
  {人はどんな状況でも常に最適な選択ができる}%
  {正解は (a)。限定合理性（bounded rationality）はサイモンの中心概念で、人は完全な情報と無限の処理能力を持たないため、最適解ではなく「満足できる解」を選ぶ（満足化）という考え方です。(d) は古典的な完全合理性の仮定です。(b)(c) は誤りで、AIにも学習データや設計に由来する限界があります。}

\quizq{5}{組織の階層と情報の関係として適切なものはどれでしょうか?}%
  {経営層は情報を必要としない}%
  {情報は階層に関係なく同じ形式で使われる}%
  {\correct{経営層ほど要約された非定型の情報を扱う傾向がある}}%
  {現場ほど要約された情報だけを使う}%
  {正解は (c)。アンソニーの3階層（戦略的計画・マネジメントコントロール・業務統制）に対応して、経営層ほど要約された非定型で外部志向・将来志向の情報を、現場ほど詳細で定型的な内部情報を扱う傾向があります。(a)(b)(d) はこの関係を否定するか逆にしています。}

